% !TEX root = ../../Main.tex
\section{Agent Design: Observation, Action and Reward}
\label{Section:AgentDesign}

The environment of \Cref{Section:TradingFramework} supplies prices and charges for trades. Three more decisions are needed to turn it into a reinforcement learning problem: what the agent sees, what it is allowed to do, and what it is rewarded for. This section states these three decisions and then describes the algorithm that learns from them. Each decision is given with the reason for it. In every case the alternative was tried and worked worse, and those failures are part of the result.

\subsection{Observation Space: Indicators, Encoding and Normalisation}

The agent observes a vector of $38$ values at each hourly bar. None of these values is a raw price. A price level carries no information that the agent can generalise from. A policy learned at a price of $1.05$ has no reason to work at $1.20$, and in eleven years of history the price moves by more than that. Every feature is therefore a ratio, and the whole observation is dimensionless. Sixteen technical indicators, listed in \Cref{Tables:Observation}, give the market context. They are split into a fast group, which describes the last hours to days, and a slow group, which describes regimes that last months. This lets the agent tell a move within a trend apart from a change of trend.

% !TEX root = ../Main.tex
\begin{table}[htb!]
\centering
\footnotesize
\caption{The sixteen technical indicators forming the market context of the observation. Periods are in hourly bars. Each moving average contributes two features, its distance from price and its own differential, both in units of the average true range; each momentum reading contributes one, scaled by the movement expected over its lookback.}
\label{Tables:Observation}
\begin{tabularx}{\textwidth}{@{}llrX@{}}
\toprule
\textbf{Indicator} & \textbf{Family} & \textbf{Period} & \textbf{Role in the observation} \\
\midrule
ATR & Volatility & 14 & Volatility unit; divides every trend feature and sets the stop distance \\
\addlinespace
RV (fast) & Volatility & 16 & Short-horizon realised volatility; divides the momentum features \\
RV (slow) & Volatility & 480 & Long-horizon realised volatility; its ratio to the fast reading gives the volatility regime \\
\addlinespace
ER & Trend & 120 & Efficiency ratio; how directional recent movement has been, bounded in $[0, 1]$ \\
\addlinespace
ROC Fast & Momentum & 24 & One day \\
ROC Medium & Momentum & 120 & One week \\
ROC Slow & Momentum & 480 & One month \\
ROC Regime & Momentum & 1{,}440 & One quarter \\
ROC Epoch & Momentum & 2{,}880 & Two quarters \\
\addlinespace
SMA Fast & Trend & 24 & One day \\
SMA Medium & Trend & 120 & One week \\
SMA Slow & Trend & 480 & One month \\
SMA Regime & Trend & 1{,}440 & One quarter \\
SMA Epoch & Trend & 2{,}880 & Two quarters \\
SMA Cycle & Trend & 4{,}320 & Three quarters \\
SMA Era & Trend & 5{,}040 & Roughly one year \\
\bottomrule
\end{tabularx}
\end{table}

The indicator values are encoded as follows. Here $C$ is the bar close, $\sigma_f$ and $\sigma_s$ are the fast and slow realised volatilities, $A$ is the average true range, and $M_i$ is the $i$-th moving average, with period $N_i$:

\begin{equation}
\frac{A}{C}, \qquad \sigma_f, \qquad \ln\!\left(\frac{\sigma_f}{\sigma_s}\right), \qquad \text{ER}
\label{Equation:VolatilityFeatures}
\end{equation}

\begin{equation}
\frac{\text{ROC}(N_j)}{\sigma_f \sqrt{N_j}}, \qquad \frac{C - M_i}{A}, \qquad \frac{M_i - M_i^{-1}}{A}
\label{Equation:TrendFeatures}
\end{equation}

Each formula removes a different dependence. Dividing the average true range by the close gives volatility as a fraction of the price. The logarithm of the ratio of fast to slow volatility shows the volatility regime. It is zero when the two agree, positive when volatility is rising and negative when it is falling. Dividing a momentum reading by $\sigma_f \sqrt{N_j}$ scales it by the movement expected over that lookback under a random walk. A value near one then means a move of normal size, whatever the period. The distance from a moving average, measured in units of the average true range, says how far the price is from its trend compared with the current volatility. The differential of the same quantity gives the direction in which the trend is moving. The other features are the position in the calendar, encoded as four sine and cosine pairs, the current signed exposure as a fraction of the largest position the account could hold, and six values that describe the last bar as log returns of its open, high, low and close against the previous close, divided by $\sigma_f$.

Features that are not bounded by construction are standardised with a causal exponentially weighted z-score, with $\alpha = 1/200$. Here $x_t$ is a raw value, and $\mu_t$ and $s^2_t$ are its running mean and variance:

\begin{equation}
z_t = \frac{x_t - \mu_{t-1}}{\sqrt{s^2_{t-1} + \epsilon}}
\label{Equation:Normalisation}
\end{equation}

\begin{equation}
\mu_t = \mu_{t-1} + \alpha \delta_t, \qquad s^2_t = (1 - \alpha)\left(s^2_{t-1} + \alpha \delta_t^2\right), \qquad \delta_t = x_t - \mu_{t-1}
\label{Equation:NormalisationUpdate}
\end{equation}

The subscript $t-1$ in \Cref{Equation:Normalisation} is the most important detail of this design. A value is standardised with statistics that end at the previous step. Only after that is the value added to those statistics. A z-score computed over the whole sample would leak future information into every observation, and it would inflate every result derived from it. Features that are already bounded, such as the exposure fraction, the efficiency ratio and the calendar encodings, have a flag that skips this layer, so their natural range is kept.

\subsection{Action Space: Position Sizing and Risk}

The agent outputs a single continuous value $a \in [-1, 1]$. Its sign is the direction of the position, and its magnitude is the fraction of a reference size to hold. Direction and size are therefore set by one decision. The discrete action sets reviewed in \Cref{Chapter:RelatedWork} cannot express this. The reference size is fixed-fractional. Let $B$ be the account balance, $\rho$ the risk fraction per position, and $k$ the stop distance in multiples of the average true range. The reference volume $V_{ref}$ is the volume that loses $\rho B$ if the price moves $k \cdot A$ against it. The target volume is this reference scaled by the action:

\begin{equation}
V_{target} = a \cdot V_{ref}, \qquad V_{ref} = \frac{\rho \, B}{k \, A}
\label{Equation:PositionSizing}
\end{equation}

with $\rho = 1\%$ and $k = 1.5$ throughout this work. The agent therefore never chooses an absolute number of units. It chooses a fraction of a size that the risk rule has already set, and this size becomes smaller on its own as volatility rises.

One defect in this calculation was found during the campaign. It is reported here because it affects one of the seven results. The volume calculation leaves out the conversion between the quote currency and the account currency. The risk actually taken is therefore $\rho / p$ and not $\rho$, where $p$ is the price of the pair. For pairs quoted close to one the difference does not matter, and the six pairs quoted against the US Dollar all take about the intended $1\%$. USD/JPY was quoted near $120$ at the start of the sample, so its effective risk is $0.008\%$. This is so small that the computed volume falls below the broker's minimum, and no position can open at all. To correct for this, $\rho$ was scaled by the price. This restores the intended risk per position and does not add leverage. \Cref{Chapter:ExperimentalResults} reports the level used for each pair. The defect is a limitation of the framework, not of the method, and it is listed as future work.

\subsection{Reward Function: Log Return, Scaling and Clipping}

The reward at each step is the log return of account equity, scaled and then clipped:

\begin{equation}
r_t = \operatorname{clip}\!\left(\lambda \ln\!\frac{E_t}{E_{t-1}},\; -c,\; +c\right), \qquad \lambda = 1000, \quad c = 1
\label{Equation:Reward}
\end{equation}

The log return is a natural choice for a trading objective because it is additive over time. The per-step rewards telescope to $\ln(E_T / E_0)$, so maximising their sum maximises the final compounded wealth. The log return also does not depend on the account size. The scale factor is needed, and the reason for it comes from a measured result. Moody and Saffell describe differential risk-adjusted rewards, which accumulate an online estimate of a Sharpe or Sortino ratio. These rewards are scale-invariant: doubling every position leaves the ratio the same. An agent that optimises such a reward gets no gradient that pushes its position size up. In this environment the policy collapses to positions of about $3\%$ of the size available to it. The resulting policy is directionally sensible, but it has no financial value. A scale-sensitive reward is therefore needed. The factor $\lambda$ raises an hourly equity log return, which is usually of order $10^{-4}$, to a magnitude that the critic can distinguish from its own approximation error.

The clip is a numerical safeguard of the kind introduced with the Deep Q-Network \cite{mnih_human-level_2015}. It is applied after the encoding. Its effect must be stated clearly here, because $\lambda$ and $c$ are not independent. Together they clip the raw log return at $\pm c / \lambda = \pm 0.1\%$. On the delivered EUR/USD model the bound is reached on $38.2\%$ of the $68\,155$ hourly steps, and the median absolute scaled reward is $0.70$. So for about two steps in five the agent is told the direction of the equity move but not its magnitude. This is a real property of the objective. It is reported here so that the learned behaviour in \Cref{Chapter:ExperimentalResults} is read against the reward the agent actually received.

\subsection{Learning Algorithm: DDPG and its Adaptations}

The algorithm is the Deep Deterministic Policy Gradient of \Cref{Section:DeepReinforcementLearning}. The actor, the critic and their delayed targets are as described there, and they are trained from a replay buffer with Ornstein-Uhlenbeck exploration noise. The equations of that section are not repeated here. The rest of this section gives the six ways in which the implementation differs from the original, each with its reason. \Cref{Tables:Hyperparameters} compares every value with the published one.

% !TEX root = ../Main.tex
\begin{table}[htb!]
\centering
\footnotesize
\caption{Hyperparameters compared with the original Deep Deterministic Policy Gradient \cite{lillicrap_continuous_2015}. Three values were changed and five mechanisms were added. The rest are the published values.}
\label{Tables:Hyperparameters}
\begin{tabularx}{\textwidth}{@{}llll X@{}}
\toprule
\textbf{Parameter} & \textbf{Original} & \textbf{This work} & \textbf{Status} & \textbf{Reason for the difference} \\
\midrule
Actor hidden layers & $400 \times 300$ & $64 \times 32$ & Changed & Noisy, non-stationary series supports less capacity \\
Discount factor $\gamma$ & $0.99$ & $0.9995$ & Changed & Horizon matched to the holding period, months and not days \\
Hidden normalisation & Batch & Layer & Changed & Inference on one observation must match training on batches \\
\addlinespace
Actor learning rate & $10^{-4}$ & $10^{-4}$ & Unchanged & \\
Critic learning rate & $10^{-3}$ & $10^{-3}$ & Unchanged & \\
Soft update rate $\tau$ & $0.001$ & $0.001$ & Unchanged & \\
Batch size $N$ & $64$ & $64$ & Unchanged & \\
Replay capacity & $10^{6}$ & $10^{6}$ & Unchanged & \\
Critic weight decay & $10^{-2}$ & $10^{-2}$ & Unchanged & \\
Noise $\theta$, $\sigma$ & $0.15$, $0.2$ & $0.15$, $0.2$ & Unchanged & \\
Final layer init & $U[-3\!\times\!10^{-3}, 3\!\times\!10^{-3}]$ & same & Unchanged & \\
\addlinespace
Gradient clip & --- & $1.0$ & Added & Bounds the size of one update caused by a rare large error \\
Warmup transitions & --- & $3{,}000$ & Added & First updates drawn from a varied buffer \\
Actor penalty $\beta$ & --- & $0.001$ & Added & Prevents saturation of the output tangent and collapse to a constant action \\
Reward scale $\lambda$ & --- & $1{,}000$ & Added & Raises an hourly log return above the critic's approximation error \\
Reward clip $c$ & --- & $1$ & Added & Numerical safeguard; limits the raw log return to $\pm 0.1\%$ \\
\bottomrule
\end{tabularx}
\end{table}

The first difference is the observation and the reward, described above.

The second is the size of the networks. The original uses hidden layers of $400$ and $300$ units for low-dimensional control tasks. Here they have $64$ and $32$ units. A $38$-dimensional observation of a noisy, non-stationary series supports much less capacity than a physics simulator with a clean reward. Larger networks fitted the training window but did not generalise beyond it. The third is the discount factor, which was raised from $0.99$ to $0.9995$. With an hourly step, $\gamma = 0.99$ gives a horizon of about a hundred bars, or about four days. This is shorter than the positions this strategy is meant to hold. At $0.9995$ the horizon grows to about two thousand bars, or roughly three months, which matches the regime indicators in the observation.

The fourth is the normalisation inside the networks. The original applies batch normalisation to the hidden layers. This implementation applies layer normalisation instead. The reason is deployment. Batch normalisation behaves differently on a single sample than on a training batch, because it depends on running statistics collected during training. On a live account, the agent must act on exactly one observation at a time. Layer normalisation normalises across the features of one sample. The computation done during training is therefore the same as the computation done in deployment. The target networks also do not need a second set of running statistics that must be kept in sync.

The fifth is gradient clipping. The original uses none. Here the gradient norm of both the actor and the critic is clipped to $1$ before each optimiser step. This limits the size of a single update. It stops a rare large temporal-difference error from moving the parameters far enough to make the critic unstable. It is only a stability measure. In particular, it does not solve the collapse discussed next, which is caused by a lack of gradient, not by too much of it.

The sixth is a penalty on the actor. This is the only change that alters the objective itself, and not just the way it is optimised. It was added because the published algorithm does not train on this problem. The actor's final layer applies a hyperbolic tangent to a pre-activation $u(s)$ to keep the action bounded. The derivative of this function vanishes as $|u|$ grows. If the policy is pushed into that region, it stops responding. The action saturates at $\pm 1$, no gradient flows back, and the agent holds one constant position for the rest of training. The fix is to penalise the magnitude of the pre-activation directly. This keeps the policy in the region where the tangent does not saturate. \Cref{Figure:ActorNetwork} shows where in the network this penalty acts:

\begin{equation}
L_{actor} = -\frac{1}{N}\sum_{i} Q\!\left(s_i, \tanh u(s_i) \,\middle|\, \theta^{Q}\right) + \frac{\beta}{N}\sum_{i} u(s_i)^2
\label{Equation:ActorRegularisation}
\end{equation}

\begin{figure}[htb!]
\centering
\resizebox{\textwidth}{!}{%
\begin{tikzpicture}[
  node distance=5mm,
  layer/.style={draw, rounded corners=1.5pt, align=center, font=\scriptsize, minimum height=9mm, inner sep=3pt, text width=20mm},
  point/.style={draw, circle, inner sep=1.5pt, font=\scriptsize},
  flow/.style={-stealth, semithick},
  pen/.style={-stealth, semithick, dashed}
]
\node[layer] (s) {observation\\$s \in \mathbb{R}^{38}$};
\node[layer, right=of s] (h1) {Linear 64\\LayerNorm\\ReLU};
\node[layer, right=of h1] (h2) {Linear 32\\LayerNorm\\ReLU};
\node[layer, right=of h2] (f3) {Linear 1};
\node[point, right=6mm of f3] (u) {$u$};
\node[layer, right=6mm of u] (th) {$\tanh$};
\node[layer, right=of th] (a) {action\\$a \in [-1,1]$};
\node[font=\scriptsize, below=8mm of u, align=center] (pen) {penalty $\beta\,\overline{u^{2}}$};

\draw[flow] (s) -- (h1);
\draw[flow] (h1) -- (h2);
\draw[flow] (h2) -- (f3);
\draw[flow] (f3) -- (u);
\draw[flow] (u) -- (th);
\draw[flow] (th) -- (a);
\draw[pen] (pen) -- (u);
\end{tikzpicture}}
\caption{The actor. Layer normalisation replaces the batch normalisation of the original, so a single observation is processed exactly as a training batch is. The penalty of \Cref{Equation:ActorRegularisation} acts on the pre-activation $u$, before the tangent. This is where saturation would otherwise remove the gradient.}
\label{Figure:ActorNetwork}
\end{figure}

with $\beta = 0.001$. Setting $\beta = 0$ recovers the original actor loss exactly, so the change extends the original loss and does not replace it. Together with the scale-sensitive reward, this penalty prevents policy collapse. The two mechanisms deal with different parts of the same failure. The reward supplies a gradient that rewards larger positions, and the penalty keeps the actor in the region where that gradient can act. Everything else follows the published algorithm. The learning rates, the soft update rate, the batch size, the replay capacity, the exploration noise parameters, the critic weight decay and the layer initialisation are the values given in the original work. A warmup period of $3000$ transitions is used before learning begins. This way the first updates are drawn from a buffer with some variety in it, and not from the first few correlated steps of an episode.